% Generated draft: requires mathematical and source review.
\documentclass[11pt]{article}
\usepackage[margin=25mm]{geometry}
\usepackage{fontspec}
\setmainfont{Latin Modern Roman}
\usepackage{amsmath,amssymb,mathtools,unicode-math}
\setmathfont{Latin Modern Math}
\usepackage{enumitem,needspace,xurl,hyperref}
\setlength{\parindent}{0pt}
\setlength{\parskip}{5pt}
\setlength{\emergencystretch}{3em}
\newcommand{\N}{\mathbb{N}}
\newcommand{\Z}{\mathbb{Z}}
\newcommand{\Q}{\mathbb{Q}}
\newcommand{\R}{\mathbb{R}}
\newcommand{\C}{\mathbb{C}}
\newcommand{\F}{\mathbb{F}}
\newcommand{\A}{\mathbb{A}}
\newcommand{\PP}{\mathbb{P}}
\newcommand{\E}{\mathbb{E}}
\newcommand{\eps}{\varepsilon}
\newcommand{\dd}{\,\mathrm{d}}
\begin{document}
\section*{UnsolvedMath ID: 30000252}
\textbf{Source problem code:} catalogue entry
\subsection{Definitions and mathematical statement}
Let $m\ge1$ and $n>2m$ be integers, and let $\Omega\subset\R^n$ be a bounded conformally contractible domain in the sense used in the report's Theorem 5. Set
\[
 p_* =\frac{n+2m}{n-2m}.
\]
The verified target asks whether its small-parameter uniqueness result for the polyharmonic problem
\[
 (-\Delta)^m u=\lambda\bigl(1+|u|^{p-1}u\bigr)\quad\text{in }\Omega,
 \qquad D^\alpha u|_{\partial\Omega}=0\quad(0\le|\alpha|\le m-1)
\]
extends from the supercritical range $p>p_*$ to the endpoint $p=p_*$. The boundary condition denotes all traces of derivatives up to order $m-1$, as in the source; it is not replaced by Navier conditions on iterated Laplacians.

More explicitly, the desired endpoint assertion retains the solution class and other hypotheses of Theorem 5 and asks whether uniqueness still holds for every sufficiently small admissible $\lambda$. In the usual positive-parameter version this means an interval $0<\lambda<\lambda_0$ for some $\lambda_0>0$ depending on the domain and fixed parameters. It must not be weakened to uniqueness only inside an arbitrarily small neighborhood of a selected solution branch.

For comparison, the standard weak Dirichlet energy space is $H_0^m(\Omega)$. At $p=p_*$ the exponent $p+1=2n/(n-2m)$ is the critical Sobolev exponent, and the weak equation pairs $(-\Delta)^m u$ with test functions in that space. This explains the endpoint terminology; the extract does not state whether Theorem 5 uses all weak solutions, classical solutions, positive solutions, or another class.

The full definition of conformal contractibility, boundary smoothness, parameter-sign convention and complete Theorem 5 are not supplied. Conformal contractibility is retained as a geometric source hypothesis and is not reduced to ordinary topological contractibility. The literature assessment checked on 21 August 2026 says no proof of the requested critical-case extension was located and cautions that the extraction is too context-dependent for a stronger classification. That dated search outcome is not independent verification of current openness.
\subsection{Short English statement}
Does the report's small-$\lambda$ uniqueness theorem for a supercritical polyharmonic Dirichlet problem remain valid at the exact critical exponent? The same domain geometry, boundary conditions and intended solution class must be retained.
\subsection{Sources}
\begin{itemize}
\item \emph{Topological and Variational Methods for Differential Equations} (2006), Oberwolfach Reports, DOI 10.4171/owr/2005/29, pp. 1601--1678, question extending Theorem 5, \url{https://doi.org/10.4171/owr/2005/29}; proposed year supplied as 2005.
\item The sole literature source in the assessment checked 21 August 2026 is the ``original problem report,'' \url{https://doi.org/10.4171/owr/2005/29}. No later endpoint theorem is supplied.
\item Local UnsolvedMath registry, permanent ID $30000252$, exact source code \texttt{catalogue entry}, created 25 August 2026 and updated 16 September 2026. Supplied verification restores the equation, domain assumptions, boundary conditions and exact critical endpoint.
\end{itemize}
\end{document}
